% DBMS Lab Project Report Template
\documentclass[12pt,a4paper]{report}
\usepackage[utf8]{inputenc}
\usepackage{graphicx}
\usepackage{hyperref}
\usepackage{geometry}
\usepackage{longtable}
\usepackage{caption}
\usepackage{float}
\usepackage{booktabs}
\usepackage{parskip}
\geometry{margin=1in}

\begin{document}

% Title Page
\begin{titlepage}
  \centering
  {\scshape\LARGE Project Title \par}
  \vspace{1.5cm}
  {\huge\bfseries Subtitle (if any)\par}
  \vspace{1.5cm}
  {\Large Team Members:\par}
  \vspace{0.5cm}
  {Name 1 --- ID1\\
   Name 2 --- ID2\\
   Name 3 --- ID3\\}
  \vfill
  Video demo: \href{https://youtu.be/your-demo-link}{https://youtu.be/your-demo-link}\\
  Code (GitHub): \href{https://github.com/your/repo}{https://github.com/your/repo}
  \vfill
  {\large Date: \today}\
\end{titlepage}

\tableofcontents
\listoffigures
\listoftables
\clearpage

% Main Sections (follow the requested structure)
\chapter{Introduction / Overview}
A brief overview of the project and its objectives. Replace this with your project's summary.

\chapter{Motivation}
Explain why you selected this project and its importance.

\chapter{Related or Similar Projects}
Discuss existing or comparable systems and cite them if needed.

\chapter{Benchmark Analysis}
Compare your implementation against existing solutions or explain benchmark methodology and results.

\chapter{Complete Feature List}
List and briefly describe every implemented feature.

\chapter{Database Design Approach}
Describe methodology (requirements analysis, normalization, indexing, transactions, backup/restore plan).

\chapter{Entity–Relationship Diagram (ERD)}
% Place your ERD image (PNG/PDF) in the folder and update the filename below.
\begin{figure}[H]
  \centering
  \includegraphics[width=0.9\textwidth]{images/erd.png}
  \caption{Entity–Relationship Diagram}
  \label{fig:erd}
\end{figure}

\chapter{Schema Diagram}
Include schema diagrams (tables, PK/FK, data types). Example:
\begin{figure}[H]
  \centering
  \includegraphics[width=0.9\textwidth]{images/schema.png}
  \caption{Database Schema Diagram}
  \label{fig:schema}
\end{figure}

\chapter{Queries for Feature Implementation}
Explain the queries you implemented for each feature. Provide grouped SQL examples.

\section*{Example: User registration (SQL)}
\begin{verbatim}
-- Create table
CREATE TABLE Users (
  Id INT PRIMARY KEY IDENTITY(1,1),
  Username VARCHAR(100) NOT NULL,
  Email VARCHAR(255) NOT NULL UNIQUE,
  PasswordHash VARCHAR(512) NOT NULL,
  CreatedAt DATETIME DEFAULT GETUTCDATE()
);

-- Example feature query: get active competitions
SELECT c.Id, c.Title, c.StartDate, c.EndDate
FROM Competitions c
WHERE c.IsApproved = 1
  AND GETUTCDATE() BETWEEN c.StartDate AND c.EndDate;
\end{verbatim}

% Point to the separate SQL file with more queries
More queries and examples are provided in the included file: \texttt{sql_queries.sql}.

\chapter{Application Screenshots}
Include screenshots of your app UI. Place images in `images/` and reference them here.

\chapter{Limitations}
List any known constraints or shortcomings.

\chapter{Future Work}
Possible improvements or extensions.

\chapter{Conclusion}
Summarize the project and key outcomes.

\chapter*{Submission Instructions}
\begin{itemize}
  \item Only one team member should submit the final PDF.
  \item The report must be prepared using LaTeX or MS Word and submitted as PDF.
  \item Ensure video link and GitHub link are live.
\end{itemize}

\vfill
\noindent{\small Generated template — replace placeholders with your content.}

\end{document}
